\documentclass[preview]{standalone}

\usepackage{amsmath}
\usepackage{amssymb}
\usepackage{stellar}
\usepackage{definitions}

\begin{document}

\id{endomorphisms-polynomial-modules}
\genpage

\section{The polynomial-module viewpoint}

\begin{snippet}{endomorphism-polynomial-module-setting}
    Let \(K\) be a \field, let \(V\) be a finite-dimensional
    \(K\)-\vectorspace, and let
    \[
        \varphi\colon V\fromto V
    \]
    be a \linearendomorphism. For \(k\in\naturalnumbers\), write
    \[
        \varphi^k
        =
        \underbrace{\varphi\circ\cdots\circ\varphi}_{k\text{ times}},
        \qquad
        \varphi^0=\identity_V.
    \]
    A choice of basis identifies \(\varphi\) with a matrix in
    \(\matrices_n(K)\), but the module viewpoint below does not depend
    on that choice.
\end{snippet}

\begin{snippetdefinition}{endomorphism-polynomial-module-definition}{Polynomial module associated with an endomorphism}
    Let \(K\) be a \field, let \(V\) be a \(K\)-\vectorspace, and let
    \(\varphi\colon V\fromto V\) be a \linearendomorphism. Define an
    action of \(K[X]\) on \(V\) by
    \[
        X\cdot v=\varphi(v)
    \]
    and, more generally, by
    \[
        \left(a_0+a_1X+\cdots+a_mX^m\right)\cdot v
        =
        a_0v+a_1\varphi(v)+\cdots+a_m\varphi^m(v).
    \]
    With this action, \(V\) is a \(K[X]\)-\module, called the
    \emph{polynomial module associated with \(\varphi\)}.
\end{snippetdefinition}

\begin{snippetproof}{endomorphism-polynomial-module-definition-proof}{endomorphism-polynomial-module-definition}{Polynomial module associated with an endomorphism}
    Additivity in the polynomial and in the vector follows from the
    linearity of \(\varphi\). If
    \[
        f=\sum_i a_iX^i,
        \qquad
        g=\sum_j b_jX^j,
    \]
    then
    \begin{align*}
        (fg)\cdot v
        &=
        \sum_{i,j}a_ib_j\varphi^{i+j}(v)\\
        &=
        \sum_i a_i\varphi^i
        \left(\sum_j b_j\varphi^j(v)\right)\\
        &=
        f\cdot(g\cdot v).
    \end{align*}
    Finally, \(1\cdot v=v\). Thus all module axioms hold.
\end{snippetproof}

\section{Invariant subspaces}

\begin{snippetdefinition}{endomorphism-invariant-subspace-definition}{Invariant subspace}
    Let \(\varphi\colon V\fromto V\) be a \linearendomorphism. A vector
    subspace \(W\leq V\) is \emph{\(\varphi\)-invariant} if
    \[
        \varphi(w)\in W
        \qquad
        \text{for every }w\in W.
    \]
\end{snippetdefinition}

\begin{snippetproposition}{invariant-subspaces-polynomial-submodules}{Invariant subspaces are polynomial submodules}
    Let \(V\) carry the \(K[X]\)-module structure associated with
    \(\varphi\). A vector subspace \(W\leq V\) is
    \(\varphi\)-\invariant \ifandonlyif it is a \(K[X]\)-\submodule of
    \(V\).
\end{snippetproposition}

\begin{snippetproof}{invariant-subspaces-polynomial-submodules-proof}{invariant-subspaces-polynomial-submodules}{Invariant subspaces are polynomial submodules}
    If \(W\) is a \(K[X]\)-\submodule, then
    \[
        \varphi(w)=X\cdot w\in W
    \]
    for every \(w\in W\), so \(W\) is invariant.

    Conversely, if \(W\) is invariant, induction gives
    \(\varphi^k(w)\in W\) for every \(k\geq0\). Hence, for
    \(f=\sum_{k=0}^{m}a_kX^k\),
    \[
        f\cdot w=\sum_{k=0}^{m}a_k\varphi^k(w)\in W.
    \]
    Thus \(W\) is closed under the \(K[X]\)-action and is a submodule.
\end{snippetproof}

\begin{snippetexample}{invariant-subspace-block-upper-triangular-example}{Invariant subspaces and block upper-triangular matrices}
    Let
    \[
        \mathcal B=\{w_1,\ldots,w_m,u_1,\ldots,u_s\}
    \]
    be a basis of \(V\), and put
    \[
        W=\operatorname{Span}_K\{w_1,\ldots,w_m\},
        \qquad
        U=\operatorname{Span}_K\{u_1,\ldots,u_s\}.
    \]
    If the matrix of \(\varphi\) in \(\mathcal B\) is
    \[
        \begin{pmatrix}
            B_1&B_2\\
            0&B_3
        \end{pmatrix},
    \]
    then \(W\) is \(\varphi\)-invariant: the images of its basis vectors
    involve only the \(w_i\). In general \(U\) is not invariant, because
    \(B_2\) records components of \(\varphi(u_j)\) along \(W\).
\end{snippetexample}

\begin{snippetproposition}{invariant-direct-sum-block-diagonal-matrix}{Invariant direct sums give block-diagonal matrices}
    Suppose
    \[
        V=W_1\oplus\cdots\oplus W_r
    \]
    and every \(W_i\) is \(\varphi\)-\invariant. Concatenating bases of
    the \(W_i\) gives a basis of \(V\) in which
    \[
        [\varphi]
        =
        \begin{pmatrix}
            B_1&0&\cdots&0\\
            0&B_2&\cdots&0\\
            \vdots&\vdots&\ddots&\vdots\\
            0&0&\cdots&B_r
        \end{pmatrix},
    \]
    where \(B_i\) represents \(\varphi|_{W_i}\). Consequently,
    \[
        \det[\varphi]=\prod_{i=1}^{r}\det B_i,
        \qquad
        \mrank[\varphi]=\sum_{i=1}^{r}\mrank B_i.
    \]
\end{snippetproposition}

\begin{snippetproof}{invariant-direct-sum-block-diagonal-matrix-proof}{invariant-direct-sum-block-diagonal-matrix}{Invariant direct sums give block-diagonal matrices}
    Since \(\varphi(W_i)\subseteq W_i\), the coordinate vector of the
    image of a basis element of \(W_i\) has no component in any
    \(W_j\) with \(j\neq i\). Hence all off-diagonal blocks vanish.
    The determinant and rank formulas are the standard formulas for a
    block-diagonal matrix.
\end{snippetproof}

\end{document}
